\section{Evaluation Metrics and Protocol}
\label{sec:evaluation}

The evaluation must show that the workflow becomes cheaper without making the agent worse. Therefore, cost metrics and quality metrics must be reported together.

\subsection{Token and Cache Metrics}

For each model call, the collector should record:
\begin{itemize}[leftmargin=2em]
    \item $\Ttotal$: total input tokens.
    \item $\Tcached$: cached input tokens reported by the provider.
    \item $\Tuncached$: uncached paid input tokens, computed as $\Ttotal - \Tcached$.
    \item $\Tout$: output tokens.
    \item $\mathrm{HitRate} = \Tcached / \Ttotal$.
\end{itemize}

The most important metric is not total input tokens. \toolname\ may leave total input unchanged. The target is lower uncached paid input:
\begin{equation}
    \mathrm{UncachedReduction}
    =
    1 - \frac{\Tuncached^{\mathrm{cache\mbox{-}friendly}}}{\Tuncached^{\mathrm{baseline}}}.
\end{equation}

\subsection{Estimated Cost}

Let $\Pin$ be uncached input price, $\Pcache$ be cached input price, and $\Pout$ be output price. The estimated actual cost of one call is:
\begin{equation}
    \cost_{\mathrm{actual}}
    =
    \Tuncached \Pin
    + \Tcached \Pcache
    + \Tout \Pout.
\end{equation}
The no-cache counterfactual is:
\begin{equation}
    \cost_{\mathrm{full}}
    =
    \Ttotal \Pin
    + \Tout \Pout.
\end{equation}
The estimated savings ratio is:
\begin{equation}
    \mathrm{Savings}
    =
    1 - \frac{\cost_{\mathrm{actual}}}{\cost_{\mathrm{full}}}.
\end{equation}

\subsection{Latency Metrics}

Prompt-cache hits should reduce repeated prefill work. We therefore track:
\begin{itemize}[leftmargin=2em]
    \item time to first token;
    \item total wall-clock latency;
    \item model-call latency excluding local tool execution when available;
    \item end-to-end task time including tool calls and tests.
\end{itemize}

\subsection{Agent Quality Metrics}

Cost improvements are only useful when task quality remains stable. For coding tasks, the primary quality metrics should be:
\begin{itemize}[leftmargin=2em]
    \item task success rate;
    \item validation pass rate, such as \texttt{cargo test}, \texttt{pytest}, or \texttt{npm test};
    \item number of turns to completion;
    \item number of tool calls;
    \item patch scope, measured by intended versus unintended files changed;
    \item reviewer score for correctness, minimality, and regressions.
\end{itemize}

\subsection{A/B Protocol}

Each benchmark task should run under two conditions:
\begin{enumerate}[leftmargin=2em]
    \item \textbf{Baseline.} A normal setup without enforced cache-friendly invariants.
    \item \textbf{Cache-friendly.} Stable provider, stable model, stable reasoning effort, Responses transport, WebSocket/session continuity where available, and no prompt rewriting.
\end{enumerate}
Both conditions should use the same task prompt, repository snapshot, model family, validation command, and success criteria. Each task should be repeated multiple times to reduce run-to-run variance.

\subsection{Benchmark Task Families}

The benchmark should include short and long sessions:
\begin{table}[h]
\centering
\caption{Suggested benchmark task families.}
\label{tab:tasks}
\small
\begin{tabularx}{\linewidth}{@{}lXX@{}}
\toprule
\textbf{Task family} & \textbf{Example} & \textbf{Why it matters} \\
\midrule
Documentation edit & Add a metric section to README & Low-risk task; exposes repeated context overhead \\
Small bug fix & Fix one failing unit test & Measures correctness under constrained edits \\
Single-file feature & Add one CLI flag and tests & Captures typical local coding work \\
Multi-file feature & Add collector, schema, and docs & Stresses longer context and tool schemas \\
Long iterative session & Improve docs across 4--6 turns & Best case for prompt-cache reuse \\
\bottomrule
\end{tabularx}
\end{table}

\subsection{Result Table Template}

\cref{tab:results-template} shows the minimum reporting table. The table is a schema, not a final result table: final values must be regenerated from the experiment JSONL paths. A positive result requires lower uncached input and no task-success regression in the all-runs accounting. Successful-only rows may be reported as diagnostics, but they must not replace the all-runs table.

\begin{table}[h]
\centering
\caption{Minimum result schema for comparing baseline and cache-friendly runs.}
\label{tab:results-template}
\small
\begin{tabularx}{\linewidth}{@{}lXX@{}}
\toprule
\textbf{Metric} & \textbf{Baseline} & \textbf{Cache-friendly} \\
\midrule
Prompt cache hit rate & from JSONL & from JSONL \\
Uncached input tokens & from JSONL & from JSONL, plus ratio to baseline \\
Estimated cost & optional price-derived value & optional price-derived value \\
Median time to first token & from JSONL when available & from JSONL when available \\
Task success, all runs & successes / records & successes / records \\
Task success, successful-only subset & diagnostic only & diagnostic only \\
\bottomrule
\end{tabularx}
\end{table}

\subsection{Current V2 Prefix Diagnostic}

The latest repository snapshot includes a fixed bounded direct-JSON diagnostic on the \texttt{docs-token-accounting} task, using Claude Code as the studied harness and \texttt{mimo-v2.5-pro} as the backend route. The run used three measured repeats for each condition on the \texttt{dynamic-drift} slice. Warm-up calls are excluded from the JSONL summary below.

\begin{table}[h]
\centering
\caption{Fixed V2 dynamic-drift diagnostic. Lower uncached input is better.}
\label{tab:v2-pilot}
\small
\begin{tabularx}{\linewidth}{@{}lrrrrr@{}}
\toprule
\textbf{Slice} & \textbf{Baseline uncached} & \textbf{Cache-friendly uncached} & \textbf{Ratio} & \textbf{Cost ratio} & \textbf{Success} \\
\midrule
\texttt{dynamic-drift} & 30{,}082 & 7{,}817 & 0.260$\times$ & 0.704$\times$ & 3/3 vs. 3/3 \\
\bottomrule
\end{tabularx}
\end{table}

The fixed diagnostic supports the narrow prefix-cache claim: moving dynamic harness state later increased cache hit rate from 91.66\% to 97.67\% and reduced paid uncached input from 30{,}082 to 7{,}817 tokens while preserving validation and task success. Observed cost also fell, but only to 0.704$\times$ baseline because output tokens increased slightly from 2{,}054 to 2{,}224. This distinction matters. The prefix intervention is incremental; it reduces paid repeated input, but it does not by itself optimize tool-output verbosity or guarantee fewer agent turns.

The earlier V2 direct-JSON pilot is retained as regression evidence rather than removed. In that pilot, the \texttt{dynamic-drift} aggregate regressed because one cache-friendly measured run used 14 Claude Code turns, compared with 6 turns for the matched baseline, increasing both repeated input and output. The runner also exposed a protocol flaw: the V2 fixture was nested under \texttt{runs/} without an isolated Git repository, so dynamic Git status could discover the parent worktree rather than only fixture-local state. After adding absolute prompt paths, \texttt{GIT\_CEILING\_DIRECTORIES}, and fixture-local Git initialization, the three-repeat diagnostic above returned to the expected direction. Because direct Claude JSON does not expose request-shape artifacts, this evidence is cache-accounting evidence rather than a complete prompt-layout trace.

\subsection{Report Regeneration}

The repository provides a report generator that turns the same JSONL files used by \texttt{eval} and \texttt{task-report} into paper-facing Markdown tables:

\begin{lstlisting}[language=bash,basicstyle=\ttfamily\small,breaklines=true,frame=single]
cargo run --quiet -- analysis-report \
  --baseline runs/<experiment>/baseline.jsonl \
  --candidate runs/<experiment>/cache-friendly.jsonl \
  --output runs/<experiment>/analysis-report.md
\end{lstlisting}

The generated report records the role separation used by the paper: Codex is the development assistant, Claude Code is the studied harness, MiMo is the current backend route/model family, and \toolname\ audit/eval logs are the measurement layer. It also emits quality, cache-accounting, and savings gates so that a lower uncached-token count is not reported as a win when task quality regresses or cache accounting is unavailable.
